% ============================================================================
% cabecalho-rodape.tex
% ----------------------------------------------------------------------------
% Define os ESTILOS DE PÁGINA (cabeçalho e rodapé) usados no documento,
% usando o pacote "fancyhdr" (carregado em includes.tex).
%
% Este arquivo cria dois estilos de página:
%   - "corpo"   -> usado nas páginas normais de conteúdo do relatório
%   - "sumario" -> usado especificamente na página do sumário (índice)
%
% Você normalmente NÃO precisa editar este arquivo. As opções de exibição
% (mostrar CNPJ, posição da numeração) são configuradas em "edicao.tex"
% através dos comandos \NumeracaoRodapetrue/false e \CNPJnoRodapetrue/false.
% ============================================================================

% -- Linhas coloridas do cabeçalho e do rodapé --------------------------------
% Usam a cor "verde-ej" definida em config.tex.
\newcommand{\linhaCabecalhoColorida}{{\color{verde-ej}\hrule width\headwidth height 1.5pt}}
\newcommand{\linhaRodapeColorida}{{\color{verde-ej}\hrule width\headwidth height 1.5pt}}

% -- Flags (chaves liga/desliga) ----------------------------------------------
% Essas flags são definidas aqui, mas seu VALOR (true/false) é escolhido em
% edicao.tex, de acordo com o comportamento desejado para cada documento.
\newif\ifCNPJnoRodape       % controla se o CNPJ aparece no rodapé
\newif\ifNumeracaoRodape    % controla se a numeração de página fica no rodapé ou no cabeçalho

% -- Numeração de página -------------------------------------------------------
% Dependendo da flag \NumeracaoRodape, coloca o número da página no
% cabeçalho (canto superior direito) ou no rodapé (canto inferior direito).
\newcommand{\numeracaoPagina}{
    \ifNumeracaoRodape
        \fancyhead[R]{}
        \fancyfoot[R]{\fontsize{10}{12}\selectfont \raisebox{18pt}[0pt][0pt]{\thepage}}
    \else
        \fancyhead[R]{\fontsize{10}{12}\selectfont \thepage}
        \fancyfoot[R]{}
    \fi
}

% -- Texto central do rodapé ---------------------------------------------------
% Monta a linha de informações institucionais que aparece no centro do
% rodapé de cada página do corpo do documento.
% Se \CNPJnoRodape estiver ativa, inclui a razão social e o CNPJ da empresa;
% caso contrário, mostra apenas razão social e endereço.
\newcommand{\rodapeCentro}{
    \ifCNPJnoRodape
        \color{gray}\fontsize{7.5}{10}\selectfont
        \razaoSocialEmpresa\ \separadorRodape\ CNPJ: \cnpjEmpresa\ \\ \empresaEndereco \separadorRodape\ \emailEmpresa\ \separadorRodape\ \telefoneEmpresa\ \\ \siteEmpresa
    \else
        \color{gray}\fontsize{7.5}{10}\selectfont
        \razaoSocialEmpresa\ \separadorRodape\ \empresaEndereco \\ \emailEmpresa\ \separadorRodape\ \telefoneEmpresa\ \separadorRodape\ \siteEmpresa
    \fi
}

% -- Estilo de página "corpo" ---------------------------------------------------
% Usado em todas as páginas normais de conteúdo (ativado em main.tex com
% \pagestyle{corpo} logo após o sumário).
\fancypagestyle{corpo}{
    \fancyhf{}  % limpa qualquer configuração de cabeçalho/rodapé anterior

    % Linhas horizontais coloridas acima do cabeçalho e abaixo do rodapé
    \renewcommand{\headrulewidth}{1pt}
    \renewcommand{\footrulewidth}{1pt}
    \renewcommand{\headrule}{\linhaCabecalhoColorida}
    \renewcommand{\footrule}{\linhaRodapeColorida}

    % Logo no canto superior esquerdo do cabeçalho.
    % Se o arquivo de logo não existir, mostra o texto "(logo)" no lugar
    % (útil para testar o modelo antes de inserir as imagens reais).
    \fancyhead[L]{
        \IfFileExists{\logoCabecalhoArquivo}{
            \includegraphics[height=\alturaLogoCabecalho]{\logoCabecalhoArquivo}
        }{
            \small(logo)
        }
    }

    \numeracaoPagina  % posiciona a numeração conforme a flag definida em edicao.tex

    \fancyfoot[C]{\rodapeCentro}  % texto institucional no centro do rodapé
}

% -- Estilo de página "sumario" --------------------------------------------------
% Usado apenas na página do sumário/índice (\thispagestyle{sumario} em
% sumario.tex). É parecido com "corpo", mas sem numeração de página.
\fancypagestyle{sumario}{
    \fancyhf{}
    \renewcommand{\headrulewidth}{1pt}
    \renewcommand{\footrulewidth}{1pt}
    \renewcommand{\headrule}{\linhaCabecalhoColorida}
    \renewcommand{\footrule}{\linhaRodapeColorida}

    \fancyhead[L]{
        \IfFileExists{\logoCabecalhoArquivo}{
            \includegraphics[height=\alturaLogoCabecalho]{\logoCabecalhoArquivo}
        }{
            \small(logo)
        }
    }
    \fancyhead[R]{}  % sem numeração no cabeçalho da página do sumário
    \fancyfoot[R]{}  % sem numeração no rodapé da página do sumário

    \fancyfoot[C]{\rodapeCentro}
}